% Figure: the value function f*(c) versus the constraint level c, with the
% Lagrange multiplier as its slope. Illustrated on the area-maximisation
% f* = P^2/16 whose slope is P/8 = lambda*. The tangent line at one budget
% shows lambda as the local marginal value of relaxing the constraint.
\begin{figure}[htbp]
\centering
\begin{tikzpicture}[scale=1.0,
  curve/.style={engblue, very thick},
  tang/.style={engred, thick},
  guide/.style={enggray, dashed},
]
  \begin{axis}[
    width=9.2cm, height=6.4cm,
    xlabel={constraint level $c$ (perimeter $P$)},
    ylabel={optimal value $f^{*}(c)$},
    xmin=0, xmax=16, ymin=0, ymax=17,
    axis lines=left,
    xtick={0,4,8,12,16}, ytick={0,4,8,12,16},
    tick label style={font=\small},
    label style={font=\small},
    clip=false,
  ]
    % Value function f* = c^2/16 (with c = P, f* = P^2/16).
    \addplot[curve, domain=0:16, samples=60] {x*x/16};
    \node[engblue, anchor=south east] at (axis cs:15.5, 15.0) {$f^{*} = c^{2}/16$};
    % Tangent at c = 8: f*(8) = 4, slope = 8/8 = 1. Line y = 4 + 1*(x-8).
    \addplot[tang, domain=2:14, samples=2] {4 + 1*(x-8)};
    \node[engred, anchor=north west] at (axis cs:12.2, 8.0)
      {slope $=\lambda^{*}=\dfrac{c}{8}$};
    % Marker at the tangent point.
    \draw[guide] (axis cs:8,0) -- (axis cs:8,4);
    \draw[guide] (axis cs:0,4) -- (axis cs:8,4);
    \filldraw[engred] (axis cs:8,4) circle (2.4pt);
  \end{axis}
\end{tikzpicture}
\caption[The value function and the multiplier as its slope.]{The value
function $f^{*}(c)$, the best achievable objective when the constraint is
$g = c$, for the area-maximisation whose optimum is $f^{*} = c^{2}/16$
(a rectangle of fixed perimeter $c = P$). The Lagrange multiplier is the
slope of this curve, $\lambda^{*} = \dv{f^{*}}{c} = c/8$, drawn as the
red tangent at $c = 8$. The multiplier reads off the marginal value of
one more unit of the constrained resource: at this operating point one
extra unit of perimeter buys $\lambda^{*} = 1$ unit of area. The rising
slope shows that the shadow price itself grows with the budget for this
problem, so the marginal value of relaxation is not constant.}
\label{fig:vol02:ch11:shadow-price}
\end{figure}
